\documentclass[../../main.tex]{subfiles}
\graphicspath{{\subfix{../../image/}}}

\begin{document}

\section{极小元的存在性}\label{section:ch08-2-极小元的存在性}

本节识别 Lagrangian $L$ 上的某些条件，以保证泛函 $I[\cdot]$ 确实存在极小元，至少在适当的 Sobolev 空间中如此。

\subsection{强制性，下半连续性}\label{subsection:ch08-2-1-强制性下半连续}

我们先就泛函
\begin{align}
I[w]\triangleq\int_U L(Dw(x),w(x),x)\,\mathrm{d}x,\label{align:ch08-2-eq01}
\end{align}
（定义于满足
\begin{align}
w=g\quad\text{在 }\partial U\text{ 上}\label{align:ch08-2-eq02}
\end{align}
的适当函数 $w:U\to\mathbb{R}$ 上）何时应有极小元，给出一些大体上启发式的洞见。

\subsubsection*{a. 强制性}

首先注意，即使是一个光滑函数 $f:\mathbb{R}\to\mathbb{R}$、下方有界，也未必取到其下确界。例如考虑 $f=e^x$ 或 $(1+x^2)^{-1}$。这些例子提示，一般而言需要某种假设来控制"大"函数 $w$ 对应的 $I[w]$。当然，确保这一点最有效的方式是假设当 $|w|\to\infty$ 时 $I[w]$"增长很快"。

更具体地，固定
\begin{align}
1<q<\infty.\label{align:ch08-2-eq03}
\end{align}
我们进一步假设
\begin{align}
\text{存在常数 }\alpha>0,\ \beta\geqslant 0\text{ 使得 }L(p,z,x)\geqslant\alpha|p|^q-\beta\quad\text{对所有 }p\in\mathbb{R}^n,\ z\in\mathbb{R},\ x\in U.\label{align:ch08-2-eq04}
\end{align}
因此
\begin{align}
I[w]\geqslant\alpha\|Dw\|_{L^q(U)}^q-\gamma,\label{align:ch08-2-eq05}
\end{align}
其中 $\gamma\triangleq\beta|U|$。于是当 $\|Dw\|_{L^q}\to\infty$ 时 $I[w]\to\infty$。习惯上称\eqref{align:ch08-2-eq05} 为 $I[\cdot]$ 的\textbf{强制性条件}。

回到寻找泛函 $I[\cdot]$ 极小元这一基本任务，由不等式\eqref{align:ch08-2-eq05} 可见，不仅对光滑函数 $w$、而且对迹意义下满足边界条件\eqref{align:ch08-2-eq02} 的 Sobolev 空间 $W^{1,q}(U)$ 中的函数 $w$ 定义 $I[w]$，似乎是合理的。毕竟，$I[w]$ 所定义的函数类 $w$ 越宽，极小元的候选者就越多。

今后写
\begin{align*}
\mathcal{A}\triangleq\{w\in W^{1,q}(U)\mid w=g\text{ 在迹意义下在 }\partial U\text{ 上}\}
\end{align*}
表示这一容许函数类。注意由\eqref{align:ch08-2-eq04}，对每个 $w\in\mathcal{A}$，$I[w]$ 均有定义（但可能等于 $+\infty$）。

\subsubsection*{b. 下半连续性}

其次注意，虽然满足强制性条件的连续函数 $f:\mathbb{R}\to\mathbb{R}$ 确实取到其下确界，但我们的积分泛函 $I[\cdot]$ 一般不会如此。为理解问题，令
\begin{align}
m\triangleq\inf_{w\in\mathcal{A}}I[w],\label{align:ch08-2-eq06}
\end{align}
并选取函数 $u_k\in\mathcal{A}$（$k=1,\cdots$）使
\begin{align}
I[u_k]\to m\quad\text{当 }k\to\infty.\label{align:ch08-2-eq07}
\end{align}
称 $\{u_k\}_{k=1}^\infty$ 为\textbf{极小化序列}。

现希望证明 $\{u_k\}_{k=1}^\infty$ 的某个子列收敛到真正的极小元。但为此需要某种紧性，而这绝对是个问题，因为空间 $W^{1,q}(U)$ 是无穷维的。事实上，若利用强制性不等式\eqref{align:ch08-2-eq05}，结果（见 \S8.2.2）只能推出极小化序列位于 $W^{1,q}(U)$ 的一个有界子集中。但这并不意味着存在某个子列在 $W^{1,q}(U)$ 中收敛。

于是我们转向弱拓扑（见 \S D.4）。由于假设 $1<q<\infty$，从而 $L^q(U)$ 自反，推得存在子列 $\{u_{k_j}\}_{j=1}^\infty\subset\{u_k\}_{k=1}^\infty$ 及函数 $u\in W^{1,q}(U)$ 使
\begin{align}
\begin{cases}
u_{k_j}\rightharpoonup u&\text{弱收敛于 }L^q(U),\\
Du_{k_j}\rightharpoonup Du&\text{弱收敛于 }L^q(U;\mathbb{R}^n).
\end{cases}\label{align:ch08-2-eq08}
\end{align}
今后把\eqref{align:ch08-2-eq08} 简写为
\begin{align}
u_{k_j}\rightharpoonup u\quad\text{弱收敛于 }W^{1,q}(U).\label{align:ch08-2-eq09}
\end{align}
此外，$u=g$ 在迹意义下在 $\partial U$ 上，故 $u\in\mathcal{A}$。

因此，通过转向弱拓扑，我们从强制性不等式\eqref{align:ch08-2-eq05} 中重新获得了足够的紧性，可对适当子列推出\eqref{align:ch08-2-eq09}。但现在出现了另一个困难：在几乎所有有意义的情形中，泛函 $I[\cdot]$ 关于弱收敛并不连续。换言之，不能由\eqref{align:ch08-2-eq07}、\eqref{align:ch08-2-eq09} 推出
\begin{align}
I[u]=\lim_{j\to\infty}I[u_{k_j}],\label{align:ch08-2-eq10}
\end{align}
从而 $u$ 是极小元。问题在于 $Du_{k_j}\rightharpoonup Du$ 并不蕴涵 $Du_{k_j}\to Du$ a.e.；例如完全可能，梯度 $Du_{k_j}$ 虽在 $L^q$ 中有界，却随 $k_j\to\infty$ 而越来越剧烈地振荡。

拯救我们的是最后一个关键观察：实际上并不需要\eqref{align:ch08-2-eq10} 的全部强度。只需知道
\begin{align}
I[u]\leqslant\liminf_{j\to\infty}I[u_{k_j}]\label{align:ch08-2-eq11}
\end{align}
即可。于是由\eqref{align:ch08-2-eq07} 可推得 $I[u]\leqslant m$。但由\eqref{align:ch08-2-eq06}，$m\leqslant I[u]$。故 $u$ 确为极小元。

\begin{definition}\label{definition:ch08-2-弱下半连续}
称泛函 $I[\cdot]$ 在 $W^{1,q}(U)$ 上是（序列）弱下半连续的，只要每当
\begin{align*}
u_k\rightharpoonup u\quad\text{弱收敛于 }W^{1,q}(U)
\end{align*}
时，均有
\begin{align*}
I[u]\leqslant\liminf_{k\to\infty}I[u_k].
\end{align*}
\end{definition}

因此我们现在的目标是在非线性项 $L$ 上识别合理条件，以保证 $I[\cdot]$ 弱下半连续。

\subsection{凸性}\label{subsection:ch08-2-2-凸性}

现回顾 \S8.1.3 中的二阶变分分析，记得在那里导出了不等式
\begin{align*}
\sum_{i,j=1}^n L_{p_ip_j}(Du(x),u(x),x)\xi_i\xi_j\geqslant 0\qquad (\xi\in\mathbb{R}^n,\ x\in U),
\end{align*}
它在 $u$ 为光滑极小元时作为必要条件成立。这一不等式强烈提示，假设 $L$ 关于其第一个自变量凸是合理的。

\begin{theorem}[弱下半连续性]\label{theorem:ch08-2-1-弱下半连续}
设 $L$ 光滑、下方有界，且此外
\begin{align*}
\text{映射 }p\mapsto L(p,z,x)\text{ 对每个 }z\in\mathbb{R},\ x\in U\text{ 是凸的。}
\end{align*}
则 $I[\cdot]$ 在 $W^{1,q}(U)$ 上弱下半连续。
\begin{proof}
1. 任取序列 $\{u_k\}_{k=1}^\infty$ 满足
\begin{align}
u_k\rightharpoonup u\quad\text{弱收敛于 }W^{1,q}(U),\label{align:ch08-2-eq12}
\end{align}
并令 $l\triangleq\liminf_{k\to\infty}I[u_k]$。需证
\begin{align}
I[u]\leqslant l.\label{align:ch08-2-eq13}
\end{align}
2. 首先由\eqref{align:ch08-2-eq12} 与 \S D.4 注意到
\begin{align}
\sup_k\|u_k\|_{W^{1,q}(U)}<\infty.\label{align:ch08-2-eq14}
\end{align}
必要时取子列，不妨还设
\begin{align}
l=\lim_{k\to\infty}I[u_k].\label{align:ch08-2-eq15}
\end{align}
此外由 \S5.7 的紧性定理可见 $u_k\to u$ 强收敛于 $L^q(U)$；必要时再取子列，有
\begin{align}
u_k\to u\quad\text{a.e. 在 }U\text{ 中}.\label{align:ch08-2-eq16}
\end{align}
3. 固定 $\varepsilon>0$。由\eqref{align:ch08-2-eq16} 与 Egoroff 定理（\S E.2），
\begin{align}
u_k\to u\quad\text{在 }E_\varepsilon\text{ 上一致},\label{align:ch08-2-eq17}
\end{align}
其中 $E_\varepsilon$ 为可测集，满足
\begin{align}
|U\setminus E_\varepsilon|\leqslant\varepsilon.\label{align:ch08-2-eq18}
\end{align}
可设当 $0<\varepsilon'<\varepsilon$ 时 $E_\varepsilon\subseteq E_{\varepsilon'}$。现写
\begin{align}
F_\varepsilon\triangleq\left\{x\in U\mid |u(x)|+|Du(x)|\leqslant\frac1\varepsilon\right\}.\label{align:ch08-2-eq19}
\end{align}
则
\begin{align}
|U\setminus F_\varepsilon|\to 0\quad\text{当 }\varepsilon\to 0.\label{align:ch08-2-eq20}
\end{align}
最后令
\begin{align}
G_\varepsilon\triangleq E_\varepsilon\cap F_\varepsilon,\label{align:ch08-2-eq21}
\end{align}
并由\eqref{align:ch08-2-eq18}、\eqref{align:ch08-2-eq20} 注意到当 $\varepsilon\to 0$ 时 $|U\setminus G_\varepsilon|\to 0$。
4. 现注意由于 $L$ 下方有界，不妨设
\begin{align}
L\geqslant 0\label{align:ch08-2-eq22}
\end{align}
（否则可对适当常数 $\beta$ 把下列论证应用于 $\widetilde{L}=L+\beta\geqslant 0$）。因此
\begin{align}
I[u_k]&=\int_U L(Du_k,u_k,x)\,\mathrm{d}x\geqslant\int_{G_\varepsilon}L(Du_k,u_k,x)\,\mathrm{d}x\notag\\
&\geqslant\int_{G_\varepsilon}L(Du,u_k,x)\,\mathrm{d}x+\int_{G_\varepsilon}D_pL(Du,u_k,x)\cdot(Du_k-Du)\,\mathrm{d}x,\label{align:ch08-2-eq23}
\end{align}
最后一不等式来自 $L$ 关于其第一个自变量的凸性；见 \S B.1。由\eqref{align:ch08-2-eq17}、\eqref{align:ch08-2-eq19} 与\eqref{align:ch08-2-eq21}，
\begin{align}
\lim_{k\to\infty}\int_{G_\varepsilon}L(Du,u_k,x)\,\mathrm{d}x=\int_{G_\varepsilon}L(Du,u,x)\,\mathrm{d}x.\label{align:ch08-2-eq24}
\end{align}
此外，由于 $D_pL(Du,u_k,x)$ 在 $G_\varepsilon$ 上一致收敛于 $D_pL(Du,u,x)$，且 $Du_k\rightharpoonup Du$ 弱收敛于 $L^q(U;\mathbb{R}^n)$，有
\begin{align}
\lim_{k\to\infty}\int_{G_\varepsilon}D_pL(Du,u_k,x)\cdot(Du_k-Du)\,\mathrm{d}x=0.\label{align:ch08-2-eq25}
\end{align}
由\eqref{align:ch08-2-eq24}、\eqref{align:ch08-2-eq25}，从\eqref{align:ch08-2-eq23} 推得
\begin{align*}
l=\lim_{k\to\infty}I[u_k]\geqslant\int_{G_\varepsilon}L(Du,u,x)\,\mathrm{d}x.
\end{align*}
此不等式对每个 $\varepsilon>0$ 成立。现令 $\varepsilon\to 0$，回顾\eqref{align:ch08-2-eq22} 与单调收敛定理（\S E.3），得
\begin{align*}
l\geqslant\int_U L(Du,u,x)\,\mathrm{d}x=I[u],
\end{align*}
即所证。
\end{proof}
\end{theorem}

\begin{remark}\label{remark:ch08-2-1-弱收敛处理}
理解上述证明如何处理弱收敛 $Du_k\rightharpoonup Du$ 非常重要。关键在于凸性不等式\eqref{align:ch08-2-eq23}，其右端 $Du_k$ 线性出现。弱收敛按其定义与线性表达式相容；故极限\eqref{align:ch08-2-eq25} 成立。请记住，一般而言 $Du_k\rightharpoonup Du$ 并不蕴涵 $Du_k\to Du$ a.e.，即使取子列亦然。
$u_k$ 到 $u$ 在 $L^q$ 中的收敛要强得多，因此我们不需要关于 $z\mapsto L(p,z,x)$ 的任何凸性假设。
\end{remark}

我们终于可以确立 $I[\cdot]$ 在 $\mathcal{A}$ 中函数之间存在极小元。

\begin{theorem}[极小元的存在性]\label{theorem:ch06-2-2-存在唯一性}
设 $L$ 满足强制性不等式\eqref{align:ch08-2-eq04} 且关于变量 $p$ 凸。又设容许集 $\mathcal{A}$ 非空。
则存在至少一个函数 $u\in\mathcal{A}$ 满足
\begin{align*}
I[u]=\min_{w\in\mathcal{A}}I[w].
\end{align*}
\begin{proof}
1. 令 $m\triangleq\inf_{w\in\mathcal{A}}I[w]$。若 $m=+\infty$，证毕；故今后设 $m$ 有限。选取极小化序列 $\{u_k\}_{k=1}^\infty$。则
\begin{align}
I[u_k]\to m.\label{align:ch08-2-eq26}
\end{align}
2. 不妨在不等式\eqref{align:ch08-2-eq04} 中取 $\beta=0$，否则同样可考虑 $\widetilde{L}=L+\beta$。于是 $L\geqslant\alpha|p|^q$，故
\begin{align}
I[w]\geqslant\alpha\int_U|Dw|^q\,\mathrm{d}x.\label{align:ch08-2-eq27}
\end{align}
由于 $m$ 有限，由\eqref{align:ch08-2-eq26}、\eqref{align:ch08-2-eq27} 得
\begin{align}
\sup_k\|Du_k\|_{L^q(U)}<\infty.\label{align:ch08-2-eq28}
\end{align}
3. 现固定任一函数 $w\in\mathcal{A}$。由于 $u_k$ 与 $w$ 均在迹意义下在 $\partial U$ 上等于 $g$，故 $u_k-w\in W_0^{1,q}(U)$。于是 Poincaré 不等式蕴涵
\begin{align*}
\|u_k\|_{L^q(U)}&\leqslant\|u_k-w\|_{L^q(U)}+\|w\|_{L^q(U)}\\
&\leqslant C\|Du_k-Dw\|_{L^q(U)}+C\leqslant C,
\end{align*}
其中最后一不等式由\eqref{align:ch08-2-eq28}。故 $\sup_k\|u_k\|_{L^q(U)}<\infty$。此估计与\eqref{align:ch08-2-eq28} 蕴涵 $\{u_k\}_{k=1}^\infty$ 在 $W^{1,q}(U)$ 中有界。
4. 因此存在子列 $\{u_{k_j}\}_{j=1}^\infty\subset\{u_k\}_{k=1}^\infty$ 及函数 $u\in W^{1,q}(U)$ 使
\begin{align*}
u_{k_j}\rightharpoonup u\quad\text{弱收敛于 }W^{1,q}(U).
\end{align*}
接着断言 $u\in\mathcal{A}$。为此注意，对如上 $w\in\mathcal{A}$，$u_k-w\in W_0^{1,q}(U)$。而 $W_0^{1,q}(U)$ 是 $W^{1,q}(U)$ 的闭线性子空间，故由 Mazur 定理（\S D.4）弱闭。因此 $u-w\in W_0^{1,q}(U)$。从而 $u$ 在 $\partial U$ 上的迹为 $g$。
于是由\cref{theorem:ch08-2-1-弱下半连续}，$I[u]\leqslant\liminf_{j\to\infty}I[u_{k_j}]=m$。但由于 $u\in\mathcal{A}$，有
\begin{align*}
I[u]=m=\min_{w\in\mathcal{A}}I[w].
\end{align*}
\end{proof}
\end{theorem}

接着转向唯一性问题。一般而言可能有多个极小元，故为保证唯一性需要进一步假设。例如假设
\begin{align}
L=L(p,x)\text{ 不依赖于 }z\label{align:ch08-2-eq29}
\end{align}
且
\begin{align}
\text{存在 }\theta>0\text{ 使 }\sum_{i,j=1}^n L_{p_ip_j}(p,x)\xi_i\xi_j\geqslant\theta|\xi|^2\quad (p,\xi\in\mathbb{R}^n;\ x\in U).\label{align:ch08-2-eq30}
\end{align}
条件\eqref{align:ch08-2-eq30} 是说对每个 $x$，映射 $p\mapsto L(p,x)$ 一致凸。

\begin{theorem}[极小元的唯一性]\label{theorem:ch08-2-3-极小元唯一}
假设\eqref{align:ch08-2-eq29}、\eqref{align:ch08-2-eq30} 成立。则 $I[\cdot]$ 的极小元 $u\in\mathcal{A}$ 在 $I[u]<\infty$ 时唯一。
\begin{proof}
1. 设 $u,\widetilde{u}\in\mathcal{A}$ 均为 $I[\cdot]$ 在 $\mathcal{A}$ 上的极小元。则 $v\triangleq\frac{u+\widetilde{u}}2\in\mathcal{A}$。断言
\begin{align}
I[v]\leqslant\frac{I[u]+I[\widetilde{u}]}2,\label{align:ch08-2-eq31}
\end{align}
且除非 $u=\widetilde{u}$ a.e.，严格不等。
2. 为此注意由一致凸性假设有
\begin{align}
L(p,x)\geqslant L(q,x)+D_pL(q,x)\cdot(p-q)+\frac{\theta}{2}|p-q|^2\qquad (x\in U,\ p,q\in\mathbb{R}^n).\label{align:ch08-2-eq32}
\end{align}
令 $q=\frac{Du+D\widetilde{u}}2$，$p=Du$，在 $U$ 上积分：
\begin{align}
I[v]+\int_U D_pL\!\left(\frac{Du+D\widetilde{u}}2,x\right)\cdot\left(\frac{Du-D\widetilde{u}}2\right)\mathrm{d}x+\frac{\theta}{8}\int_U|Du-D\widetilde{u}|^2\,\mathrm{d}x\leqslant I[u].\label{align:ch08-2-eq33}
\end{align}
类似地，在\eqref{align:ch08-2-eq32} 中令 $q=\frac{Du+D\widetilde{u}}2$，$p=D\widetilde{u}$ 并积分：
\begin{align}
I[v]+\int_U D_pL\!\left(\frac{Du+D\widetilde{u}}2,x\right)\cdot\left(\frac{D\widetilde{u}-Du}2\right)\mathrm{d}x+\frac{\theta}{8}\int_U|Du-D\widetilde{u}|^2\,\mathrm{d}x\leqslant I[\widetilde{u}].\label{align:ch08-2-eq34}
\end{align}
两式相加除以 2，得
\begin{align*}
I[v]+\frac{\theta}{8}\int_U|Du-D\widetilde{u}|^2\,\mathrm{d}x\leqslant\frac{I[u]+I[\widetilde{u}]}2.
\end{align*}
此即证明\eqref{align:ch08-2-eq31}。
3. 由于 $I[u]=I[\widetilde{u}]=\min_{w\in\mathcal{A}}I[w]\leqslant I[v]$，推得 $Du=D\widetilde{u}$ a.e. 在 $U$ 中。又 $u=\widetilde{u}=g$ 在迹意义下在 $\partial U$ 上，故 $u=\widetilde{u}$ a.e.
\end{proof}
\end{theorem}

\subsection{Euler--Lagrange 方程的弱解}\label{subsection:ch08-2-3-EulerLagrange弱解}

接下来希望证明 $I[\cdot]$ 的任一极小元 $u\in\mathcal{A}$ 在某种适当意义下求解 Euler--Lagrange 方程。这不能由 \S8.1 的计算推出，因为我们不知道 $u$ 光滑，只知道 $u\in W^{1,q}(U)$。事实上还需要 $L$ 及其导数上的一些增长条件。今后假设
\begin{align}
|L(p,z,x)|\leqslant C(|p|^q+|z|^q+1),\label{align:ch08-2-eq35}
\end{align}
以及
\begin{align}
\begin{cases}
|D_pL(p,z,x)|\leqslant C(|p|^{q-1}+|z|^{q-1}+1),\\
|L_z(p,z,x)|\leqslant C(|p|^{q-1}+|z|^{q-1}+1),
\end{cases}\label{align:ch08-2-eq36}
\end{align}
其中 $C$ 为常数，对所有 $p\in\mathbb{R}^n$、$z\in\mathbb{R}$、$x\in U$ 成立。

\paragraph{弱解定义的动机。}现转向与泛函 $L$ 相关联的 Euler--Lagrange PDE 的边值问题，对光滑极小元 $u$ 写为
\begin{align}
\begin{cases}
-\sum_{i=1}^n\big(L_{p_i}(Du,u,x)\big)_{x_i}+L_z(Du,u,x)=0&\text{在 }U\text{ 中},\\
u=g&\text{在 }\partial U\text{ 上}.
\end{cases}\label{align:ch08-2-eq37}
\end{align}
若以检验函数 $v\in C_c^\infty(U)$ 乘\eqref{align:ch08-2-eq37} 并分部积分，得到等式
\begin{align}
\int_U\sum_{i=1}^n L_{p_i}(Du,u,x)v_{x_i}+L_z(Du,u,x)v\,\mathrm{d}x=0.\label{align:ch08-2-eq38}
\end{align}
这当然就是 \S8.1.2 推导\eqref{align:ch08-2-eq37} 时首次得到的恒等式。

现设 $u\in W^{1,q}(U)$。利用\eqref{align:ch08-2-eq36} 可见
\begin{align*}
|D_pL(Du,u,x)|\leqslant C(|Du|^{q-1}+|u|^{q-1}+1)\in L^{q'}(U),
\end{align*}
其中 $q'=\frac{q}{q-1}$，$\frac1q+\frac1{q'}=1$。类似地
\begin{align}
|L_z(Du,u,x)|\leqslant C(|Du|^{q-1}+|u|^{q-1}+1)\in L^{q'}(U).\label{align:ch08-2-eq39}
\end{align}
因此用标准逼近论证可知，等式\eqref{align:ch08-2-eq38} 对任一 $v\in W_0^{1,q}(U)$ 均成立。这引出下列定义。

\begin{definition}\label{definition:ch08-2-EulerLagrange弱解}
称 $u\in\mathcal{A}$ 为边值问题\eqref{align:ch08-2-eq37}（Euler--Lagrange 方程）的弱解，只要
\begin{align*}
\int_U\sum_{i=1}^n L_{p_i}(Du,u,x)v_{x_i}+L_z(Du,u,x)v\,\mathrm{d}x=0
\end{align*}
对所有 $v\in W_0^{1,q}(U)$ 成立。
\end{definition}

\begin{theorem}[Euler--Lagrange 方程的解]\label{theorem:ch08-2-4-EulerLagrange弱解}
设 $L$ 满足增长条件\eqref{align:ch08-2-eq35}、\eqref{align:ch08-2-eq36}，且 $u\in\mathcal{A}$ 满足
\begin{align*}
I[u]=\min_{w\in\mathcal{A}}I[w].
\end{align*}
则 $u$ 是\eqref{align:ch08-2-eq37} 的弱解。
\begin{proof}
如 \S8.1.2 那样进行，但需小心在积分号下求导。固定任一 $v\in W_0^{1,q}(U)$，令
\begin{align*}
i(\tau)\triangleq I[u+\tau v]\qquad (\tau\in\mathbb{R}).
\end{align*}
由\eqref{align:ch08-2-eq35} 知对所有 $\tau$，$i(\tau)$ 有限。

取 $\tau\neq 0$，写差商
\begin{align}
\frac{i(\tau)-i(0)}{\tau}&=\int_U\frac{L(Du+\tau Dv,u+\tau v,x)-L(Du,u,x)}{\tau}\,\mathrm{d}x\notag\\
&=\int_U L_\tau(x)\,\mathrm{d}x,\label{align:ch08-2-eq40}
\end{align}
其中
\begin{align*}
L_\tau(x)\triangleq\frac1\tau\big[L(Du(x)+\tau Dv(x),u(x)+\tau v(x),x)-L(Du(x),u(x),x)\big]
\end{align*}
对 a.e. $x\in U$。显然
\begin{align}
L_\tau(x)\to\sum_{i=1}^n L_{p_i}(Du,u,x)v_{x_i}+L_z(Du,u,x)v\quad\text{a.e.}\label{align:ch08-2-eq41}
\end{align}
当 $\tau\to 0$。此外
\begin{align*}
L_\tau(x)&=\frac1\tau\int_0^\tau\frac{\mathrm{d}}{\mathrm{d}s}L(Du+sDv,u+sv,x)\,\mathrm{d}s\\
&=\frac1\tau\int_0^\tau\sum_{i=1}^n L_{p_i}(Du+sDv,u+sv,x)v_{x_i}+L_z(Du+sDv,u+sv,x)v\,\mathrm{d}s.
\end{align*}
回顾 \S B.2 的 Young 不等式：$ab\leqslant\frac{a^q}{q}+\frac{b^{q'}}{q'}$，其中 $\frac1q+\frac1{q'}=1$。由于 $u,v\in W^{1,q}(U)$，不等式\eqref{align:ch08-2-eq36} 与 Young 不等式经初等计算后给出
\begin{align*}
|L_\tau(x)|\leqslant C(|Du|^q+|u|^q+|Dv|^q+|v|^q+1)\in L^1(U)
\end{align*}
对每个 $\tau\neq 0$。因此可应用控制收敛定理，由\eqref{align:ch08-2-eq40}、\eqref{align:ch08-2-eq41} 得 $i'(0)$ 存在且等于
\begin{align*}
\int_U\sum_{i=1}^n L_{p_i}(Du,u,x)v_{x_i}+L_z(Du,u,x)v\,\mathrm{d}x.
\end{align*}
但由于 $i(\cdot)$ 在 $\tau=0$ 处取最小值，知 $i'(0)=0$；故 $u$ 是弱解。
\end{proof}
\end{theorem}

\begin{remark}\label{remark:ch08-2-2-凸联合弱解即极小}
一般而言，Euler--Lagrange 方程\eqref{align:ch08-2-eq37} 还有其他不对应于 $I[\cdot]$ 最小值的解；见 \S8.2.5 与 \S8.5。但在联合映射 $(p,z)\mapsto L(p,z,x)$ 对每个 $x$ 凸的特殊情形，每个弱解实际上都是极小元。
为看到这一点，设 $u\in\mathcal{A}$ 在弱意义下解
\begin{align}
\begin{cases}
-\sum_{i=1}^n\big(L_{p_i}(Du,u,x)\big)_{x_i}+L_z(Du,u,x)=0&\text{在 }U\text{ 中},\\
u=g&\text{在 }\partial U\text{ 上},
\end{cases}\label{align:ch08-2-eq42}
\end{align}
并任取 $w\in\mathcal{A}$。利用映射 $(p,z)\mapsto L(p,z,x)$ 的凸性，有
\begin{align*}
L(p,z,x)+D_pL(p,z,x)\cdot(q-p)+L_z(p,z,x)\cdot(w-z)\leqslant L(q,w,x).
\end{align*}
令 $p=Du(x)$，$q=Dw(x)$，$z=u(x)$，$w=w(x)$，在 $U$ 上积分：
\begin{align*}
I[u]+\int_U D_pL(Du,u,x)\cdot(Dw-Du)+L_z(Du,u,x)(w-u)\,\mathrm{d}x\leqslant I[w].
\end{align*}
由\eqref{align:ch08-2-eq42}，左端第二项为零，故对每个 $w\in\mathcal{A}$ 有 $I[u]\leqslant I[w]$。
\end{remark}

\subsection{方程组}\label{subsection:ch08-2-4-方程组}

\subsubsection*{a. 凸性}

现再次采用 \S8.1.4 中为方程组设定的记号，考察泛函
\begin{align*}
I[\mathbf{w}]\triangleq\int_U L(D\mathbf{w}(x),\mathbf{w}(x),x)\,\mathrm{d}x
\end{align*}
（定义于适当函数 $\mathbf{w}:U\to\mathbb{R}^m$，其中 $L:\mathbb{M}^{m\times n}\times\mathbb{R}^m\times\overline{U}\to\mathbb{R}$ 给定）的极小元存在性问题。

\S8.2.2 中建立的理论可毫无困难地推广到当前情形。故假设强制性不等式
\begin{align}
L(P,z,x)\geqslant\alpha|P|^q-\beta\qquad (P\in\mathbb{M}^{m\times n},\ z\in\mathbb{R}^m,\ x\in U)\label{align:ch08-2-eq43}
\end{align}
（$\alpha>0$，$\beta\geqslant 0$），并令
\begin{align*}
\mathcal{A}\triangleq\{\mathbf{w}\in W^{1,q}(U;\mathbb{R}^m)\mid \mathbf{w}=\mathbf{g}\text{ 在迹意义下在 }\partial U\text{ 上}\},
\end{align*}
其中 $\mathbf{g}:\partial U\to\mathbb{R}^m$ 给定。

\begin{theorem}[极小元的存在性]\label{theorem:ch08-2-5-方程组存在}
设 $L$ 满足强制性不等式\eqref{align:ch08-2-eq43} 且关于变量 $P$ 凸。又设容许集 $\mathcal{A}$ 非空。
则存在 $\mathbf{u}\in\mathcal{A}$ 满足
\begin{align*}
I[\mathbf{u}]=\min_{\mathbf{w}\in\mathcal{A}}I[\mathbf{w}].
\end{align*}
证明几乎逐字遵循 \S8.2.2 中\cref{theorem:ch08-2-1-弱下半连续} 与\cref{theorem:ch06-2-2-存在唯一性} 的证明。
\end{theorem}

与上面\cref{theorem:ch08-2-3-极小元唯一} 类似，有

\begin{theorem}[极小元的唯一性]\label{theorem:ch08-2-6-方程组唯一}
设 $L$ 不依赖于 $z$，且映射 $P\mapsto L(P,x)$ 一致凸。则 $I[\cdot]$ 的极小元 $\mathbf{u}\in\mathcal{A}$ 在 $I[\mathbf{u}]<\infty$ 时唯一。
\end{theorem}

现此外假设
\begin{align}
\begin{cases}
|L(P,z,x)|\leqslant C(|P|^q+|z|^q+1),\\
|D_PL(P,z,x)|\leqslant C(|P|^{q-1}+|z|^{q-1}+1),\\
|D_zL(P,z,x)|\leqslant C(|P|^{q-1}+|z|^{q-1}+1),
\end{cases}\label{align:ch08-2-eq44}
\end{align}
其中 $C$ 为常数，对所有 $P\in\mathbb{M}^{m\times n}$、$z\in\mathbb{R}^m$、$x\in U$ 成立。

现考虑 Euler--Lagrange 方程组
\begin{align}
\begin{cases}
-\sum_{i=1}^n\big(L_{p_i^k}(D\mathbf{u},\mathbf{u},x)\big)_{x_i}+L_{z^k}(D\mathbf{u},\mathbf{u},x)=0&\text{在 }U\text{ 中},\\
u^k=g^k&\text{在 }\partial U\text{ 上},
\end{cases}\label{align:ch08-2-eq45}
\end{align}
$k=1,\cdots,m$，并定义 $\mathbf{u}\in\mathcal{A}$ 为弱解，只要
\begin{align*}
\sum_{k=1}^m\int_U\sum_{i=1}^n L_{p_i^k}(D\mathbf{u},\mathbf{u},x)w_{x_i}^k+L_{z^k}(D\mathbf{u},\mathbf{u},x)w^k\,\mathrm{d}x=0
\end{align*}
对所有 $\mathbf{w}\in W_0^{1,q}(U;\mathbb{R}^m)$，$\mathbf{w}=(w^1,\cdots,w^m)$ 成立。

\begin{theorem}[Euler--Lagrange 方程组的解]\label{theorem:ch08-2-7-方程组弱解}
设 $L$ 满足增长条件\eqref{align:ch08-2-eq44}，且 $\mathbf{u}\in\mathcal{A}$ 满足
\begin{align*}
I[\mathbf{u}]=\min_{\mathbf{w}\in\mathcal{A}}I[\mathbf{w}].
\end{align*}
则 $\mathbf{u}$ 是\eqref{align:ch08-2-eq45} 的弱解。
证明几乎与\cref{theorem:ch08-2-4-EulerLagrange弱解} 完全相同。
\end{theorem}

\subsubsection*{b. 多凸性}

相当出人意料的是，存在一些数学上、物理上有趣的方程组，并未被上述\cref{theorem:ch08-2-5-方程组存在} 覆盖，但仍可用变分法研究。其中包括某些问题，其 Lagrangian $L$ 关于 $P$ 并不凸，但 $I[\cdot]$ 仍弱下半连续。

\begin{lemma}[行列式的弱连续性]\label{lemma:ch08-2-行列式弱连续}
设 $n<q<\infty$，且
\begin{align*}
\mathbf{u}_k\rightharpoonup\mathbf{u}\quad\text{弱收敛于 }W^{1,q}(U;\mathbb{R}^n).
\end{align*}
则
\begin{align*}
\det D\mathbf{u}_k\rightharpoonup\det D\mathbf{u}\quad\text{弱收敛于 }L^{q/n}(U).
\end{align*}
\begin{proof}
1. 首先回顾矩阵恒等式 $(\det P)I=P(\operatorname{cof}P)^T$；故
\begin{align*}
\det P=\sum_{j=1}^n p_j^i(\operatorname{cof}P)_j^i\qquad (i=1,\cdots,n).
\end{align*}
2. 现令 $\mathbf{w}\in C^\infty(U;\mathbb{R}^n)$，$\mathbf{w}=(w^1,\cdots,w^n)$。则
\begin{align}
\det D\mathbf{w}=\sum_{j=1}^n w_{x_j}^i(\operatorname{cof}D\mathbf{w})_j^i\qquad (i=1,\cdots,n).\label{align:ch08-2-eq46}
\end{align}
但 \S8.1.4 的引理断言 $\sum_{j=1}^n(\operatorname{cof}D\mathbf{w})_{x_j}^i{}_j=0$。故公式\eqref{align:ch08-2-eq46} 表明
\begin{align*}
\det D\mathbf{w}=\sum_{j=1}^n\big(w^i(\operatorname{cof}D\mathbf{w})_j^i\big)_{x_j}.
\end{align*}
因此梯度矩阵的行列式可写成散度。于是若 $v\in C_c^\infty(U)$，有
\begin{align}
\int_U v\det D\mathbf{w}\,\mathrm{d}x=-\sum_{j=1}^n\int_U v_{x_j}w^i(\operatorname{cof}D\mathbf{w})_j^i\,\mathrm{d}x\qquad (i=1,\cdots,n).\label{align:ch08-2-eq47}
\end{align}
3. 我们已对光滑函数 $\mathbf{w}$ 建立恒等式\eqref{align:ch08-2-eq47}，标准逼近论证给出
\begin{align}
\int_U v\det D\mathbf{u}_k\,\mathrm{d}x=-\sum_{j=1}^n\int_U v_{x_j}u_k^i(\operatorname{cof}D\mathbf{u}_k)_j^i\,\mathrm{d}x\label{align:ch08-2-eq48}
\end{align}
（$k=1,2,\cdots$）。由于 $n<q<\infty$ 且 $\mathbf{u}_k\rightharpoonup\mathbf{u}$ 于 $W^{1,q}(U;\mathbb{R}^n)$，由 Morrey 不等式知 $\{\mathbf{u}_k\}_{k=1}^\infty$ 在 $C^{0,1-n/q}(U;\mathbb{R}^n)$ 中有界。故用 Arzelà--Ascoli 紧性准则（\S C.8）推得 $\mathbf{u}_k\to\mathbf{u}$ 在 $U$ 上一致。回到恒等式\eqref{align:ch08-2-eq48}，可见若知道
\begin{align}
\lim_{k\to\infty}\int_U\psi(\operatorname{cof}D\mathbf{u}_k)_j^i\,\mathrm{d}x=\int_U\psi(\operatorname{cof}D\mathbf{u})_j^i\,\mathrm{d}x\label{align:ch08-2-eq50}
\end{align}
（对 $i,j=1,\cdots,n$ 及每个 $\psi\in C_c^\infty(U)$），则可得出
\begin{align}
\lim_{k\to\infty}\int_U v\det D\mathbf{u}_k\,\mathrm{d}x=-\sum_{j=1}^n\int_U v_{x_j}u^i(\operatorname{cof}D\mathbf{u})_j^i\,\mathrm{d}x=\int_U v\det D\mathbf{u}\,\mathrm{d}x.\label{align:ch08-2-eq49}
\end{align}
然而 $(\operatorname{cof}D\mathbf{u}_k)_j^i$ 是一个 $(n-1)\times(n-1)$ 矩阵的行列式，可按上述方式分析：写成适当 $(n-2)\times(n-2)$ 子矩阵行列式之和，乘以一致收敛因子。继续下去，最终只需证明显然事实：矩阵 $D\mathbf{u}_k$ 的各元弱收敛于 $D\mathbf{u}$ 的相应元。由此验证\eqref{align:ch08-2-eq50} 从而\eqref{align:ch08-2-eq49}。
4. 最后，由于 $\{\mathbf{u}_k\}_{k=1}^\infty$ 在 $W^{1,q}(U;\mathbb{R}^n)$ 中有界且 $|\det D\mathbf{u}_k|\leqslant C|D\mathbf{u}_k|^n$，知 $\{\det D\mathbf{u}_k\}_{k=1}^\infty$ 在 $L^{q/n}(U)$ 中有界。故任一子列在 $L^{q/n}(U)$ 中有弱收敛子列——由\eqref{align:ch08-2-eq49}——只能收敛到 $\det D\mathbf{u}$。
\end{proof}
\end{lemma}

现利用此引理建立与\cref{theorem:ch08-2-1-弱下半连续} 类似的弱下半连续断言，但不假设 Lagrangian $L$ 关于 $P$ 必凸。反之设 $m=n$ 且 $L$ 形如
\begin{align}
L(P,z,x)=F(P,\det P,z,x)\qquad (P\in\mathbb{M}^{n\times n},\ z\in\mathbb{R}^n,\ x\in U),\label{align:ch08-2-eq51}
\end{align}
其中 $F:\mathbb{M}^{n\times n}\times\mathbb{R}\times\mathbb{R}^n\times\overline{U}\to\mathbb{R}$ 光滑。进一步假设
\begin{align}
\text{对每个固定的 }z\in\mathbb{R}^m,\ x\in\mathbb{R}^n,\text{ 联合映射 }(P,r)\mapsto F(P,r,z,x)\text{ 凸。}\label{align:ch08-2-eq52}
\end{align}
形如\eqref{align:ch08-2-eq51}、满足\eqref{align:ch08-2-eq52} 的 Lagrangian $L$ 称为\textbf{多凸的}。

\begin{theorem}[多凸泛函的下半连续性]\label{theorem:ch08-2-8-多凸下半连续}
设 $n<q<\infty$。又设 $L$ 下方有界且多凸。则 $I[\cdot]$ 在 $W^{1,q}(U;\mathbb{R}^n)$ 上弱下半连续。
\begin{proof}
任取序列 $\{\mathbf{u}_k\}_{k=1}^\infty$ 满足
\begin{align}
\mathbf{u}_k\rightharpoonup\mathbf{u}\quad\text{弱收敛于 }W^{1,q}(U;\mathbb{R}^n).\label{align:ch08-2-eq53}
\end{align}
由引理，
\begin{align}
\det D\mathbf{u}_k\rightharpoonup\det D\mathbf{u}\quad\text{弱收敛于 }L^{q/n}(U).\label{align:ch08-2-eq54}
\end{align}
现可几乎完全如同\cref{theorem:ch08-2-1-弱下半连续} 的证明那样论证。事实上，用该证明的记号，有
\begin{align*}
I[\mathbf{u}_k]&=\int_U L(D\mathbf{u}_k,\mathbf{u}_k,x)\,\mathrm{d}x\geqslant\int_{G_\varepsilon}L(D\mathbf{u}_k,\mathbf{u}_k,x)\,\mathrm{d}x\\
&=\int_{G_\varepsilon}F(D\mathbf{u}_k,\det D\mathbf{u}_k,\mathbf{u}_k,x)\,\mathrm{d}x\\
&\geqslant\int_{G_\varepsilon}F(D\mathbf{u},\det D\mathbf{u},\mathbf{u}_k,x)\,\mathrm{d}x\\
&\quad+\int_{G_\varepsilon}F_P(D\mathbf{u},\det D\mathbf{u},\mathbf{u}_k,x)\cdot(D\mathbf{u}_k-D\mathbf{u})\\
&\qquad\qquad+F_r(D\mathbf{u},\det D\mathbf{u},\mathbf{u}_k,x)(\det D\mathbf{u}_k-\det D\mathbf{u})\,\mathrm{d}x,
\end{align*}
此由\eqref{align:ch08-2-eq52}。如\cref{theorem:ch08-2-1-弱下半连续} 证明中那样推理，由\eqref{align:ch08-2-eq53}、\eqref{align:ch08-2-eq54} 推得最后一项当 $k\to\infty$ 时极限为零。
\end{proof}
\end{theorem}

如前，立即推得

\begin{theorem}[多凸泛函的极小元存在性]\label{theorem:ch08-2-9-多凸存在}
设 $n<q<\infty$，$L$ 满足强制性不等式\eqref{align:ch08-2-eq43} 且多凸。又设容许集 $\mathcal{A}$ 非空。
则存在 $\mathbf{u}\in\mathcal{A}$ 满足
\begin{align*}
I[\mathbf{u}]=\min_{\mathbf{w}\in\mathcal{A}}I[\mathbf{w}].
\end{align*}
\end{theorem}

\begin{example}[弹性力学]\label{example:ch08-2-弹性}
对多凸性的大部分兴趣来自非线性弹性理论，其中 $n=3$。考虑一个弹性体，其初始参考构型为 $U$。然后把每个边界点 $x\in\partial U$ 位移到新位置 $\mathbf{g}(x)$，希望确定每个内点 $x\in U$ 的新位移 $\mathbf{u}(x)$。
若材料是超弹性的，则按定义存在一个相关联的能量密度 $L$，使物理位移 $\mathbf{u}$ 在所有容许位移 $\mathbf{w}\in\mathcal{A}$ 上极小化内能泛函
\begin{align*}
I[\mathbf{w}]\triangleq\int_U L(D\mathbf{w},x)\,\mathrm{d}x.
\end{align*}
从物理上看，表示拉伸与压缩内能密度的 $L$ 显式依赖于局部体积变化，即依赖于 $\det D\mathbf{w}$，似乎是合理的。换言之，物理上宜假设 $L$ 具有形式 $L(P,x)=F(P,\det P,x)$。于是 $F$ 在其第一个自变量中描述由线元变化引起的内能变化，在其第二个自变量中描述由体积元变化引起的内能变化。更多解释见 Ball（Arch.\ Rational Mech.\ Analysis 63 (1977), 337--403）。
\end{example}

\subsection{局部极小元}\label{subsection:ch08-2-5-局部极小元}

本小节提出一个基本问题：在何种情形下，能量泛函 $I[\cdot]$ 的临界点实际上是极小元或局部极小元？具体地，假设 $u$ 是 Euler--Lagrange PDE
\begin{align}
\begin{cases}
-\sum_{i=1}^n\big(L_{p_i}(Du,u,x)\big)_{x_i}+L_z(Du,u,x)=0&\text{在 }U\text{ 中},\\
u=g&\text{在 }\partial U\text{ 上}
\end{cases}\label{align:ch08-2-eq55}
\end{align}
的光滑解，因而是泛函
\begin{align*}
I[w]=\int_U L(Dw,w,x)\,\mathrm{d}x
\end{align*}
（在满足边界条件 $w=g$ 在 $\partial U$ 上的函数 $w$ 之间）的临界点。如常假设
\begin{align*}
p\mapsto L(p,z,x)\text{ 凸}.
\end{align*}

我们将证明：若图像 $x\mapsto u(x)$ 位于由其他临界点对应的单参数图像族 $x\mapsto u(x,\lambda)$ 生成的区域 $R$ 内，则事实上 $u$ 是 $I[\cdot]$ 在取值于 $R$ 内的容许变分 $w$ 之间的极小元。更精确地，设 $I\subseteq\mathbb{R}$ 为包含 $0$ 的开区间，且 $\{u(\cdot,\lambda)\mid\lambda\in I\}$ 是 Euler--Lagrange PDE
\begin{align}
-\sum_{i=1}^n\big(L_{p_i}(Du(x,\lambda),u(x,\lambda),x)\big)_{x_i}+L_z(Du(x,\lambda),u(x,\lambda),x)=0\label{align:ch08-2-eq56}
\end{align}
在 $U$ 内的光滑单参数解族，满足
\begin{align}
u(x)=u(x,0)\qquad (x\in U).\label{align:ch08-2-eq57}
\end{align}

如下构造一个取值于区域 $R$（即函数 $u(\cdot,\lambda)$（$\lambda\in I$）图像之并）中的容许函数 $w$。取 $\theta:\overline{U}\to I$ 为满足
\begin{align}
\theta=0\quad\text{在 }\partial U\text{ 上}\label{align:ch08-2-eq58}
\end{align}
的光滑函数。定义
\begin{align}
w(x)\triangleq u(x,\theta(x)),\label{align:ch08-2-eq59}
\end{align}
注意 $w=u=g$ 在 $\partial U$ 上。

\begin{theorem}[局部极小元]\label{theorem:ch08-2-10-局部极小元}
函数 $u$ 是区域 $R$ 内的局部极小元，意指
\begin{align}
I[u]\leqslant I[w]\label{align:ch08-2-eq60}
\end{align}
对任何如上构造的函数 $w$ 成立。
因此，若 $u$ 是 Euler--Lagrange PDE 的解并嵌入一族其他解之中，则 $u$ 是 $I[\cdot]$ 在形如\eqref{align:ch08-2-eq59} 的函数 $w$ 之间的极小元。例如，若对小 $\lambda$ 有 $u_\lambda>0$，则任一在点态意义下充分接近 $u$ 的 $w$ 均可写成此形式。注意 $Dw$ 不必接近 $Du$。
\begin{proof}
1. 首先观察
\begin{align*}
w_{x_i}(x)=u_{x_i}(x,\theta(x))+u_\lambda(x,\theta(x))\theta_{x_i}\qquad (i=1,\cdots,n).
\end{align*}
因此 $L$ 关于其第一个自变量的凸性蕴涵
\begin{align}
I[w]&=\int_U L(Dw,w,x)\,\mathrm{d}x=\int_U L(Du+u_\lambda D\theta,w,x)\,\mathrm{d}x\notag\\
&\geqslant\int_U L(Du,w,x)+u_\lambda D_pL(Du,w,x)\cdot D\theta\,\mathrm{d}x,\label{align:ch08-2-eq61}
\end{align}
其中 $u$ 在 $(x,\theta(x))$ 处取值，$D=D_x$。
2. 现引入向量场 $\mathbf{b}=(b^1,\cdots,b^n)$，定义为
\begin{align}
b^i\triangleq\int_0^{\theta(x)}u_\lambda(x,\lambda)L_{p_i}(Du(x,\lambda),u(x,\lambda),x)\,\mathrm{d}\lambda\qquad (i=1,\cdots,n).\label{align:ch08-2-eq62}
\end{align}
则
\begin{align*}
\operatorname{div}\mathbf{b}&=\sum_{i=1}^n\theta_{x_i}u_\lambda(x,\theta(x))L_{p_i}(Du(x,\theta(x)),u(x,\theta(x)),x)\\
&\quad+\sum_{i=1}^n\int_0^{\theta(x)}u_{\lambda x_i}(x,\lambda)L_{p_i}(Du(x,\lambda),u(x,\lambda),x)\\
&\qquad+u_\lambda(x,\lambda)\big(L_{p_i}(Du(x,\lambda),u(x,\lambda),x)\big)_{x_i}\,\mathrm{d}\lambda\\
&=u_\lambda(x,\theta)D_pL(Du(x,\theta(x)),w(x),x)\cdot D\theta\\
&\quad+\int_0^{\theta(x)}\sum_{i=1}^n u_{\lambda x_i}(x,\lambda)L_{p_i}(Du(x,\lambda),u(x,\lambda),x)\\
&\qquad+u_\lambda(x,\lambda)L_z(Du(x,\lambda),u(x,\lambda),x)\,\mathrm{d}\lambda.
\end{align*}
此计算中最后一等式调用了 Euler--Lagrange 方程\eqref{align:ch08-2-eq56}。接着观察
\begin{align*}
\big(L(Du(x,\lambda),u(x,\lambda),x)\big)_\lambda=\sum_{i=1}^n L_{p_i}(Du,u,x)u_{x_i\lambda}+L_z(Du,u,x)u_\lambda.
\end{align*}
因此上述计算蕴涵
\begin{align*}
\operatorname{div}\mathbf{b}&=u_\lambda(x,\theta)D_pL(Du,w,x)\cdot D\theta+\int_0^{\theta(x)}L(Du,u,x)_\lambda\,\mathrm{d}\lambda\\
&=u_\lambda(x,\theta)D_pL(Du,w,x)\cdot D\theta+L(Du(x,\theta(x)),w,x)-L(Du,u,x),
\end{align*}
此由\eqref{align:ch08-2-eq57}。
3. 把此计算用于\eqref{align:ch08-2-eq61}，用 Gauss--Green 定理推得
\begin{align*}
I[w]&\geqslant\int_U L(Du,u,x)+\operatorname{div}\mathbf{b}\,\mathrm{d}x\\
&=I[u]+\int_{\partial U}\mathbf{b}\cdot\nu\,\mathrm{d}S=I[u],
\end{align*}
因为 $\theta=0$ 在 $\partial U$ 上，从而 $\mathbf{b}=0$ 在 $\partial U$ 上。\eqref{align:ch08-2-eq60} 得证。
\end{proof}
\end{theorem}

这一计算让我们得以一瞥关于临界点何时为极小元或局部极小元的深刻经典理论：更多见 Morrey [Mo] 与 Giaquinta--Hildebrandt [G-H]。

\end{document}
